\section{Purpose}

This document defines the minimum formatting, writing, accessibility, and document-control requirements for documentaion.

These requirements are intended to ensure that documents are:
\begin{itemize}
    \item clear and easy to read;
    \item usable on screen and when printed;
    \item consistent across departments and procedures;
    \item accessible to users with different visual and reading needs;
    \item suitable for safe and repeatable laboratory work; and
    \item easy to review, approve, revise, and archive.
\end{itemize}

\section{Scope}

This standard applies to all new and substantially revised documents. This includes, but is not limited to:
\begin{itemize}
    \item standard operating procedures;
    \item laboratory test and analysis methods;
    \item sample handling procedures;
    \item equipment operating instructions;
    \item cleaning, maintenance, and calibration procedures;
    \item safety and waste-disposal procedures;
    \item quality control and quality assurance procedures;
    \item work instructions and checklists; and
    \item forms, templates, and records associated with the documents above.
\end{itemize}

This standard applies to documents distributed electronically, printed, or both. It applies to the source files, generated PDF files, and printed copies where those formats are under the control of the document owner.

This standard covers:
\begin{itemize}
    \item page size, margins, binding offsets, headers, footers, and page numbering;
    \item typography, colour, spacing, headings, lists, tables, and figures;
    \item the wording and presentation of warnings, cautions, and notices;
    \item accessibility and readability requirements;
    \item units, symbols, abbreviations, dates, and numerical values;
    \item document structure and procedure-writing conventions; and
    \item document identification, revision history, approval, and review.
\end{itemize}

This standard does not replace technical, legal, regulatory, safety, or quality-system requirements. Where another applicable requirement is more stringent than this standard, the more stringent requirement \must{} be followed. Any conflict or approved exception \must{} be documented and authorised by the responsible document owner.

Documents created before the effective date of this standard \should{} be updated during their next scheduled review. They \may{} be updated earlier when changes to the procedure, equipment, hazards, regulations, or quality system require revision.

Unless an approved exception exists, all documents within the scope of this standard \must{} conform to the requirements marked \must{}. Requirements marked \should{} are recommended and \may{} be omitted only where there is a documented reason. Requirements marked \may{} are optional. Requirements marked \mustnot{} or \shouldnot{} prohibit or discourage the stated practice, respectively.